\section{Resultados}

\subsection{Caudal de Ventilación}

Aplicando la Ec. (\ref{eq:caudal}) con un volumen efectivo de 320 m$^3$ y 12 renovaciones de aire por hora, se obtiene un caudal de ventilación de 64.0 m$^3$/min. La equivalencia en otras unidades se presenta en la Tabla \ref{tab:flow_results}.

\begin{table}[htbp]
	\centering
	\caption{Resultados del caudal de ventilación}
	\label{tab:flow_results}
	\setcounter{itemcount}{0}
	\begin{tabularx}{\textwidth}{@{}NXllc@{}}
		\toprule
		\multicolumn{1}{c}{\textbf{Item}} & \textbf{Magnitud} & \textbf{Valor} & \textbf{Unidad} & \textbf{Referencia} \\
		\midrule
		& Caudal volumétrico & 64.0 & m$^3$/min & Ec. (\ref{eq:caudal}) \\
		& Caudal volumétrico & 3 840 & m$^3$/h & Conversión \\
		& Caudal volumétrico & 2 260.1 & CFM & Conversión \\
		& Caudal en unidades SI & 1.0667 & m$^3$/s & Conversión \\
		\bottomrule
	\end{tabularx}
\end{table}

\subsection{Potencia del Ventilador}

La Tabla \ref{tab:fan_power} resume los resultados de potencia para diferentes escenarios de presión total. El escenario de diseño corresponde a 250 Pa, que incluye pérdidas por filtración, accesorios y rejillas de exfiltración.

\begin{table}[htbp]
	\centering
	\caption{Resultados de potencia del ventilador}
	\label{tab:fan_power}
	\setcounter{itemcount}{0}
	\begin{tabularx}{\textwidth}{@{}NXllc@{}}
		\toprule
		\multicolumn{1}{c}{\textbf{Item}} & \textbf{Escenario} & \textbf{Presión Total} & \textbf{Potencia} & \textbf{Referencia} \\
		& & \textbf{(Pa)} & \textbf{(kW)} & \\
		\midrule
		& Conservador & 150 & 0.267 & Ec. (\ref{eq:potencia}) \\
		& Diseño & 250 & 0.444 & Ec. (\ref{eq:potencia}) \\
		& Alto & 350 & 0.622 & Ec. (\ref{eq:potencia}) \\
		& Motor recomendado & -- & 1.0 HP & Criterio de diseño \\
		\bottomrule
	\end{tabularx}
\end{table}

\subsection{Área de Impulsión del Ventilador}

Con una velocidad de impulsión de 8 m/s, el área de boca del ventilador es 0.1333 m$^2$, lo cual corresponde a un diámetro equivalente de 412 mm y un radio de 206.0 mm. Los resultados se presentan en la Tabla \ref{tab:inlet_results}.

\begin{table}[htbp]
	\centering
	\caption{Resultados del área de impulsión del ventilador}
	\label{tab:inlet_results}
	\setcounter{itemcount}{0}
	\begin{tabularx}{\textwidth}{@{}NXllc@{}}
		\toprule
		\multicolumn{1}{c}{\textbf{Item}} & \textbf{Parámetro} & \textbf{Valor} & \textbf{Unidad} & \textbf{Referencia} \\
		\midrule
		& Área de impulsión & 0.1333 & m$^2$ & Ec. (\ref{eq:area_vent}) \\
		& Diámetro equivalente & 412 & mm & Ec. (\ref{eq:diametro}) \\
		& Radio equivalente & 206.0 & mm & Ec. (\ref{eq:diametro}) \\
		& Velocidad de impulsión & 8.0 & m/s & Criterio de diseño \\
		\bottomrule
	\end{tabularx}
\end{table}

\subsection{Rejillas de Exfiltración}

Para una velocidad de exfiltración de 3.0 m/s, el área neta total requerida es 0.3556 m$^2$. Distribuida en tres rejillas, cada una debe tener un área neta de 0.1185 m$^2$. Las dimensiones calculadas son 353.2 mm de alto por 335.5 mm de ancho. La Tabla \ref{tab:outlet_results} resume estos resultados.

\begin{table}[htbp]
	\centering
	\caption{Resultados de las rejillas de exfiltración}
	\label{tab:outlet_results}
	\setcounter{itemcount}{0}
	\begin{tabularx}{\textwidth}{@{}NXllc@{}}
		\toprule
		\multicolumn{1}{c}{\textbf{Item}} & \textbf{Parámetro} & \textbf{Valor} & \textbf{Unidad} & \textbf{Referencia} \\
		\midrule
		& Número de rejillas & 3 & unidades & Criterio de diseño \\
		& Velocidad de exfiltración & 3.0 & m/s & Criterio de diseño \\
		& Área neta total & 0.3556 & m$^2$ & Ec. (\ref{eq:area_exfil}) \\
		& Área neta por rejilla & 0.1185 & m$^2$ & División equitativa \\
		& Altura de rejilla & 353.2 & mm & Criterio de diseño \\
		& Ancho de rejilla & 335.5 & mm & Criterio de diseño \\
		\bottomrule
	\end{tabularx}
\end{table}

\subsection{Datos para Modelado CFD}

La Tabla \ref{tab:cfd_boundaries} presenta las condiciones de contorno recomendadas para el modelo CFD.

\begin{table}[htbp]
	\centering
	\caption{Condiciones de contorno para el modelo CFD}
	\label{tab:cfd_boundaries}
	\setcounter{itemcount}{0}
	\begin{tabularx}{\textwidth}{@{}Nlllc@{}}
		\toprule
		\multicolumn{1}{c}{\textbf{Item}} & \textbf{Elemento} & \textbf{Geometría} & \textbf{Dimensiones} & \textbf{Condición de Contorno} \\
		\midrule
		& Impulsión del ventilador & Círculo & $r$ = 206.0 mm & Velocity inlet, 8 m/s \\
		& Rejilla de salida 1 & Rectángulo & 353 mm $\times$ 336 mm & Velocity outlet, 3 m/s \\
		& Rejilla de salida 2 & Rectángulo & 353 mm $\times$ 336 mm & Velocity outlet, 3 m/s \\
		& Rejilla de salida 3 & Rectángulo & 353 mm $\times$ 336 mm & Velocity outlet, 3 m/s \\
		& Paredes del laboratorio & Sólidas & Según geometría 3D & Wall, no slip \\
		\bottomrule
	\end{tabularx}
\end{table}

El balance de masas entre entrada y salida presenta una discrepancia inferior al 0.1 \%. El caudal de entrada es 1.0667 m$^3$/s, mientras que el caudal de salida es 3 $\times$ (0.3532 m $\times$ 0.3355 m) $\times$ 3 m/s = 1.0665 m$^3$/s. La diferencia (0.02 \%) se debe al redondeo de dimensiones.
